\chapter{Fundamentals}
\label{ch:fundamentals}

\section{What is a legacy software development project?}
\subsection{Legacy and legacy software}
When talking about legacy in the realm of computer technology it is often used to describe a previous or outdated computer system.
Or more in general a technology from an earlier time \citep{Legacy}.

This definition also fits well with software in general, which can still be in use, but is not actively maintained or developed anymore.

\subsection{Legacy code}
Legacy code on the other hand is not simply only outdated software, but its challenges go beyond its age to the complexities of its structure and the difficulties involved in modifying it.
The core issue with legacy code is the absence of comprehensive tests, which increases the risk of introducing bugs during updates.
This creates a culture of fear among developers about changes, which leads to overly cautious development practices and discourages necessary refactoring.
To deal with legacy code, it is essential to transform it into a state that allows for safe modifications through incremental improvements.
This should be underpinned by a rigorous emphasis on testing.
Adequate testing reveals how the code works, allowing for quicker and more informed decision-making.
This demystifies the process of working with legacy code and reduces the anxieties it typically provokes \citep{feathersWorkingEffectivelyLegacy2004}.

There is no exact definition of what legacy code is, however, in general, it is characterized by the following aspects:
\begin{itemize}
    \item \textbf{Lack of Testing:} Legacy code tends to lack proper unit tests, integration tests, or automated test suites. This absence of testing makes it risky to modify or extend.
    \item \textbf{Sparse Documentation:} Documentation for legacy software is often inadequate or missing altogether. This lack of documentation makes it challenging for developers to understand the system's behavior and design.
    \item \textbf{Difficulty in Understanding:} Legacy code can be hard to comprehend, whether or not documentation exists. Its complexity, convoluted logic, and reliance on outdated programming paradigms or patterns contribute to this difficulty.
\end{itemize}
\citep{feathersWorkingEffectivelyLegacy2004}.

Following these aspects, even new code can become legacy code if it is not maintained and developed properly.

In conclusion, legacy code and the resulting lack of changing code, also results in a software being automatically legacy software over time due to technology  development projects are not using the latest technologies and development practices anymore and are therefore considered legacy software development projects.

\subsection{Legacy software development project}

There is no exact definition of what a legacy software development project is, but to summarize:

it can be described as a software development project that is actively maintained and developed, but is not using the latest technologies and development practices anymore.

In the context of software development, it is often referred to as software that is outdated, unsupported, and often difficult to maintain. This can be due to the age of the software, the technology stack it is built on, or the lack of documentation. In this thesis, the term legacy software development project is used to describe a software development project that is actively maintained and developed, but is not using the latest technologies and development practices. This can be due to the age of the project, the technology stack it is built on, or the lack of resources to update the project.

\section{Age of active software development projects}
\section{Software development life cycle}
\section{Reproducibility in software engineering}
\section{Methods to create reproducible results in software engineering}
\subsection{Scripted rollout of native developer machines}
\subsection{Virtual machines}
\subsection{Containers}
\section{Containerization in software development}
\subsection{DevContainers}
\subsection{IDE integration}

\section{Performance measurement}

\section{General definitions}

\section{Summary}
